% abstract en inglés (una sola línea \noindent, sin \par)
We present the design and implementation of a molecular dynamics (MD)
simulation pipeline for studying the inelastic rebound of Lennard-Jones (LJ)
clusters on an FCC argon crystalline surface. Rather than an \textit{ad hoc}
approach, the work emphasizes software architecture as a first-class research
component: metaclasses, descriptors, parametric decorators, generators, typed
exception hierarchies, and \texttt{pandas}/\texttt{scipy} analysis are all
integrated with the \textsc{Lammps} MD engine. The central observable is the
coefficient of restitution $e$ as a function of cluster size
$N\in\{1,2,7,13,55\}$, incident speed $v_0$, angle $\theta$, and
projectile-surface adhesion $\varepsilon_{ps}$. A total of \num{1350}
simulations were performed in a parallel parametric sweep, revealing that
$\alpha_{E,\mathrm{int}}(N)$ grows systematically with $N$ while the COR
decreases, consistent with energy dissipation into internal vibrational modes.
The Weber number $\mathrm{We}=E_{\mathrm{inc}}/(N\varepsilon_{ps})$ is
introduced as a scaling variable unifying rebound and sticking regimes. The
fully version-controlled, reproducible source code constitutes a pedagogical
case study in advanced scientific programming.
